\section{Reproducibility Statement}
\label{sec:repro}
\paragraph{Declaration.} Every number in this report comes from the code and settings in our public repository \url{https://github.com/Prakash9188/Adversarially_Robust_Face_Verification} (\texttt{configs/default.yaml}, seed 0). The thresholds, the flag threshold and the detector type were all fixed on validation data before we looked at the test split. The splits do not share any person (checked in the code and in a unit test). All results are from a single seed on CPU (Windows 11, Python 3.11.9, PyTorch 2.15.0+cpu). Because PGD uses a random start, numbers may differ slightly on other hardware or library versions.

\paragraph{What is in the repository.} Code (\texttt{src/}, \texttt{scripts/}), 17 unit tests, a pinned environment (\texttt{requirements-lock.txt}, \texttt{environment.yml}), result tables and figures (\texttt{experiments/results/}) and the source of this report (\texttt{report/}). The tables and figures in this report are generated from the result CSV files by \texttt{scripts/make\_report\_assets.py}.

\paragraph{Checkpoints.} We did not train any model. The only checkpoint is the public VGGFace2 InceptionResnetV1 from \texttt{facenet-pytorch}~\cite{esler2019facenetpytorch}, which downloads by itself on the first run (about 107\,MB); we do not share it ourselves. The detector is refit in a few minutes from the cached attacks, so we do not store it. \textbf{Data:} LFW is not in the repository. The code downloads LFW-funneled from the figshare mirror (\url{https://ndownloader.figshare.com/files/5976015}), which has the same archive as the original LFW page.

\paragraph{Steps.} Make a virtual environment and install the pinned packages (see README). Then run \texttt{python -m pytest tests}, \texttt{python scripts/run\_baseline.py}, \texttt{python scripts/run\_detector.py} (takes hours on CPU; attacks are cached) and \texttt{python scripts/make\_report\_assets.py}.

\paragraph{Commits.} We developed the project in three commits (hashes and times below, also in \texttt{docs/commit\_log.md}).

\begin{table}[h]
\centering\footnotesize
\setlength{\tabcolsep}{3pt}
\resizebox{\columnwidth}{!}{\begin{tabular}{clll}
\toprule
\# & Hash & Timestamp & Message \\
\midrule
1 & \texttt{59ff3b2} & 2026-11-26 16:26:01+05:30 & Implement project foundation and baseline methodology \\
2 & \texttt{7720d6a} & 2026-11-26 19:07:50+05:30 & Implement proposed novelty and experimental improvements \\
3 & \texttt{\hashthree} & \timethree & Finalize project, evaluation, and documentation \\
\bottomrule
\end{tabular}
}
\caption{Development commits. Commit 3 contains this report, so its hash is kept in the repository log (\texttt{docs/commit\_log.md}) and not in the PDF it ships with.}
\label{tab:commits}
\end{table}
